\section{实践案例与治理}

\subsection{两个真实案例的对比}

考虑两个典型团队的 reward learning 流程：OpenAI 在 2022 年的 InstructGPT 中大量使用 human preference data，而 Stanford RL 团队则在机器人任务里靠 goal classifiers。下表总结出两者的落差：

\begin{table}[H]
\centering
\begin{tabular}{p{0.25\textwidth} p{0.2\textwidth} p{0.2\textwidth} p{0.25\textwidth}}
\toprule
维度 & OpenAI (InstructGPT) & Stanford Robots & 结论 \\
\midrule
数据来源 & 人类 preference comparison & success/failure snapshots & 人类标签更精准但更贵 \\
Reward signal & sigmoid difference & binary classifier & Bradley-Terry 更适合对话，goal classifier 更适合 deterministic tasks \\
RL 算法 & PPO + KL & TRPO + KL flow & 都加入 KL 罚项避免 drifting \\
部署 & API + guardrails & 仿真 → 真实 robot & API 需要更多安全审查 \\
\bottomrule
\end{tabular}
\caption{两个 reward learning 流程的对比}
\end{table}

\begin{warningbox}{治理角度的教训}
在 API 端，OpenAI 通过 Rate limiting & guardrail filters 补充 reward model；在机器人端，Stanford 团队更依赖 instrumentation + simulators。因此治理策略必须和任务类型匹配。
\end{warningbox}

\subsection{Observability Playbook}

强烈建议按照以下 checklist 构建可观测性 playbook：

\begin{enumerate}
    \item 记录每条 trajectory 的 reward model score + KL penalty
    \item 将 reward anomalies 査到 dashboard（例如 sudden drop below 0.1）
    \item 自动化检测过度重复、fast reward jumps、out-of-domain inputs
    \item 每周 review failure logs，与人类偏好团队沟通是否需要更多 data
\end{enumerate}

\begin{importantbox}{自动化报警的设定}
报警规则可以基于模型 score 的 moving average：当 reward score window 上升 3 个 std 或 sudden drop 0.2+ 时推送到 PagerDuty 并 trigger rollback script。
\end{importantbox}

\subsection{本章小结}

实践案例表明 reward learning 既要匹配任务也要匹配组织治理。Checklist 与可观测性是防止 reward hacking 的关键。

